\documentclass[a4paper,11pt]{article}

\usepackage[T1]{fontenc}
\usepackage{lmodern}
\usepackage[utf8]{inputenc}
\usepackage[indonesian]{babel}
\usepackage{jheppub}
\usepackage{booktabs,longtable,array}
\usepackage{tikz}
\usetikzlibrary{shapes.geometric,arrows.meta,positioning,calc}

\renewcommand{\acknowledgments}{\section*{Ucapan Terima Kasih}}
\AtBeginDocument{\renewcommand{\proofname}{Bukti}}

\tikzset{
  blok/.style={rectangle,draw,rounded corners=2pt,align=center,text width=5.6cm,minimum height=0.75cm,font=\footnotesize,inner sep=3pt},
  kecil/.style={rectangle,draw,rounded corners=2pt,align=center,text width=3.9cm,minimum height=0.7cm,font=\footnotesize,inner sep=3pt},
  samping/.style={rectangle,draw,dashed,align=center,text width=3.4cm,minimum height=0.6cm,font=\scriptsize,inner sep=2pt},
  ujung/.style={rectangle,draw,rounded corners=8pt,align=center,minimum height=0.65cm,font=\footnotesize,inner sep=4pt},
  kep/.style={diamond,draw,aspect=2.6,align=center,text width=3.1cm,inner sep=0pt,font=\footnotesize},
  panah/.style={-{Stealth[length=2mm]}}
}

\title{\boldmath Optimasi Biaya Pengiriman Komoditas Kesehatan HIV/AIDS pada Data SCMS-USAID dengan Model Transportasi: Solusi Awal NWCR, LCM, dan VAM yang Dioptimalkan dengan MODI}

\author[a,1]{[Nama Anggota 1 (NRP)],\note{Penulis korespondensi.}}
\author[a]{[Nama Anggota 2 (NRP)],}
\author[a]{[Nama Anggota 3 (NRP)]}
\author[a]{dan [Nama Anggota 4 (NRP)]}

\affiliation[a]{[Nama Departemen], Institut Teknologi Sepuluh Nopember,\\Surabaya, Jawa Timur, Indonesia 60111}

\emailAdd{[email anggota 1]}

\subheader{Tugas Akhir Riset Operasi 1, Kelompok [nomor kelompok]}

\abstract{Artikel ini memodelkan pengiriman obat dan alat tes HIV/AIDS yang dikelola Supply Chain Management System (SCMS) milik USAID sebagai masalah transportasi. Data pengiriman 2006--2015 sebanyak 10.324 baris dibersihkan menjadi 6.037 pengiriman, yang membentuk 22 node sumber (21 kelompok produsen dan satu gudang regional SCMS), 38 node tujuan berupa negara penerima, dan 258 rute teramati. Pasokan dan permintaan diambil dari total berat pengiriman sehingga model berimbang pada 19.550.180 kg, sedangkan biaya per kilogram dihitung sebagai rata-rata tertimbang pada tiap rute. Rute yang tidak pernah muncul di data (578 dari 836 sel) dilarang dengan biaya Big-M. Solusi awal dicari dengan North-West Corner Rule (NWCR), Least Cost Method (LCM), dan Vogel's Approximation Method (VAM), kemudian masing-masing diperbaiki dengan metode MODI sampai optimal. NWCR dan LCM semula menaruh 12,3\% dan 1,4\% muatan pada rute terlarang, sedangkan VAM langsung layak dengan biaya 56,69 juta USD. Ketiganya berakhir pada biaya optimal yang sama, yaitu 54,46 juta USD, atau 19,83\% di bawah biaya aktual 67,93 juta USD. Perbedaannya hanya pada jumlah iterasi MODI: 97 untuk NWCR, 42 untuk LCM, dan 16 untuk VAM.}
\keywords{model transportasi, NWCR, LCM, VAM, MODI, rantai pasok}

\notoc

\begin{document}
\maketitle
\flushbottom

\section{Pendahuluan}
\label{sec:pendahuluan}

\subsection{Latar belakang}
Pengobatan HIV/AIDS bergantung pada pasokan obat antiretroviral (ARV) dan alat tes yang datang tepat waktu. Di banyak negara berpendapatan rendah, barang-barang itu tidak diproduksi di dalam negeri. Barang harus dibeli dari pabrik di India, Eropa, Jepang, atau Amerika, lalu dikirim lintas benua ke gudang atau klinik di negara penerima. Salah satu operator logistik terbesar untuk keperluan ini adalah Supply Chain Management System (SCMS), yang dijalankan dengan dukungan USAID dalam kerangka program PEPFAR \cite{usaid2015}. Sebagian pengiriman SCMS dikirim langsung dari pabrik ke negara tujuan (\textit{direct drop}), sebagian lain disinggahkan dulu di gudang regional (\textit{regional distribution center}, RDC) sebelum diteruskan.

Ongkos angkut bukan angka kecil. Pada sampel yang dianalisis di sini, 6.037 pengiriman menghabiskan biaya freight sekitar 67,9 juta USD untuk memindahkan hampir 19,6 ribu ton barang. Biaya per kilogram juga sangat tidak seragam: antar pengiriman, nilainya berkisar dari kurang dari 0,1 sampai lebih dari 1.600 USD/kg, bergantung pada moda, jarak, dan ukuran muatan. Karena itu pertanyaan ``barang dari sumber mana sebaiknya dikirim ke negara mana'' punya dampak finansial yang nyata.

\subsection{Alasan pemilihan model transportasi}
Struktur masalahnya cocok dengan model transportasi. Ada sejumlah titik pasokan dengan jumlah barang yang terbatas, ada sejumlah titik permintaan dengan kebutuhan tertentu, dan ongkos angkut per satuan barang berbeda untuk setiap pasangan sumber--tujuan. Yang dicari adalah jumlah barang pada tiap pasangan yang membuat total ongkos paling kecil. Jika diperiksa satu per satu, tersedia $22\times38=836$ pasangan, dan jumlah kombinasi alokasi yang memenuhi pasokan dan permintaan tidak terhingga banyaknya. Mencoba-coba secara manual jelas tidak masuk akal.

Pendekatan lain kurang pas untuk data ini. \textit{Dynamic programming} lebih cocok untuk keputusan bertahap, sedangkan model jaringan umum lebih berguna ketika ada simpul perantara dengan aturan aliran yang rumit. Data SCMS hanya mencatat asal, tujuan, berat, dan ongkos tiap pengiriman, sehingga bentuk paling alami adalah dua himpunan node (sumber dan tujuan) yang dihubungkan oleh busur berbiaya. Itulah definisi model transportasi \cite{hitchcock1941,koopmans1949}.

\subsection{Rumusan masalah, tujuan, dan batasan}
Makalah ini menjawab empat pertanyaan. Pertama, bagaimana data pengiriman riil yang berantakan diubah menjadi masalah transportasi lengkap dengan node, pasokan, permintaan, dan biaya. Kedua, bagaimana solusi awal dari NWCR, LCM, dan VAM bila diterapkan pada data tersebut. Ketiga, bagaimana solusi awal itu diperbaiki dengan MODI hingga optimal dan berapa iterasi yang dibutuhkan masing-masing. Keempat, seberapa jauh biaya optimal berbeda dari biaya aktual.

Cakupan penelitian dibatasi pada biaya freight dari sumber ke tujuan. Metode yang dibandingkan adalah NWCR, LCM, dan VAM sebagai pencari solusi awal, dan MODI sebagai pemeriksa serta pemerbaik solusi. Seluruh algoritma ditulis sendiri tanpa memakai \textit{solver} program linear.

\section{Landasan Teori}
\label{sec:teori}

\subsection{Model transportasi}
Misalkan ada $m$ sumber dengan pasokan $a_i$ dan $n$ tujuan dengan permintaan $b_j$. Ongkos mengirim satu satuan barang dari sumber $i$ ke tujuan $j$ adalah $c_{ij}$, dan jumlah yang dikirim adalah $x_{ij}$. Modelnya:
\begin{subequations}\label{eq:model}
\begin{align}
\min\; Z &= \sum_{i=1}^{m}\sum_{j=1}^{n} c_{ij}\,x_{ij}, \label{eq:obj}\\
\text{dengan}\quad \sum_{j=1}^{n} x_{ij} &= a_i, \quad i=1,\dots,m, \label{eq:sup}\\
\sum_{i=1}^{m} x_{ij} &= b_j, \quad j=1,\dots,n, \label{eq:dem}\\
x_{ij} &\ge 0 .\label{eq:nn}
\end{align}
\end{subequations}
Model dikatakan berimbang jika $\sum_i a_i=\sum_j b_j$. Pada kondisi itu, satu solusi dasar layak memuat paling banyak $m+n-1$ variabel dasar. Untuk data ini $m+n-1=22+38-1=59$.

\subsection{North-West Corner Rule (NWCR)}
NWCR memulai alokasi dari sel pojok kiri atas tanpa memperhatikan biaya. Langkahnya: (1) alokasikan $\min(a_i,b_j)$ pada sel yang sedang ditinjau; (2) kurangi pasokan dan permintaan dengan jumlah tersebut; (3) bila pasokan baris habis pindah ke baris berikutnya, bila permintaan kolom habis pindah ke kolom berikutnya; (4) ulangi sampai semua terpenuhi. Metode ini sangat cepat tetapi hasilnya sepenuhnya ditentukan oleh urutan baris dan kolom \cite{taha2017}. Dalam penelitian ini baris dan kolom diurutkan dari volume terbesar ke terkecil.

\subsection{Least Cost Method (LCM)}
LCM memilih sel dengan biaya terkecil di antara baris dan kolom yang masih aktif, mengalokasikan sebanyak mungkin, lalu mencoret baris atau kolom yang sudah terpenuhi. Proses berhenti ketika tidak ada baris atau kolom aktif. Bila ada dua sel berbiaya sama, dipilih yang bisa menampung alokasi lebih besar. Karena langsung melihat biaya, LCM umumnya lebih baik daripada NWCR, tetapi pilihannya serakah: sel murah diambil lebih dulu tanpa memikirkan akibat bagi sel lain \cite{hillier2021}.

\subsection{Vogel's Approximation Method (VAM)}
VAM menambahkan gagasan biaya peluang. Pada setiap iterasi dihitung penalti tiap baris dan kolom aktif, yaitu selisih antara dua biaya terkecil di baris atau kolom itu:
\begin{equation}
p_i=c_{i}^{(2)}-c_{i}^{(1)}, \qquad q_j=c_{j}^{(2)}-c_{j}^{(1)},
\label{eq:penalti}
\end{equation}
dengan $c^{(1)}$ dan $c^{(2)}$ masing-masing biaya terkecil pertama dan kedua. Baris atau kolom dengan penalti terbesar dipilih, lalu sel berbiaya terkecil di dalamnya diisi semaksimal mungkin. Penalti besar berarti kerugian besar bila sel termurah tidak dipakai, sehingga garis itu perlu didahulukan \cite{reinfeld1958}. Jika hanya tersisa satu elemen aktif pada suatu garis, penaltinya diambil sama dengan biayanya.

\subsection{Metode MODI}
\textit{Modified distribution} (MODI), atau metode $u$--$v$, dipakai untuk memeriksa dan memperbaiki solusi dasar layak \cite{taha2017,hillier2021}. Untuk setiap sel basis berlaku $u_i+v_j=c_{ij}$. Dengan menetapkan $u_1=0$, seluruh $u_i$ dan $v_j$ diperoleh dari $m+n-1$ persamaan tersebut. Untuk sel non-basis dihitung biaya reduksi (\textit{reduced cost})
\begin{equation}
d_{ij}=c_{ij}-u_i-v_j .
\label{eq:dij}
\end{equation}
Jika semua $d_{ij}\ge 0$, solusi sudah optimal. Jika masih ada yang negatif, sel dengan $d_{ij}$ paling negatif dipilih sebagai sel masuk. Sel masuk bersama sel-sel basis membentuk tepat satu siklus tertutup. Sel-sel pada siklus diberi tanda $+$ dan $-$ bergantian, dimulai dengan $+$ pada sel masuk. Dengan $\theta=\min x_{ij}$ pada sel bertanda $-$, nilai $\theta$ ditambahkan ke sel bertanda $+$ dan dikurangkan dari sel bertanda $-$. Sel bertanda $-$ yang menjadi nol keluar dari basis. Perubahan biaya total pada satu iterasi adalah $\theta\,d_{ij}\le 0$, sehingga biaya tidak pernah naik. Proses diulang sampai tidak ada $d_{ij}$ negatif.

\subsection{Penanganan rute yang tidak tersedia}
Tidak semua pasangan sumber--tujuan pernah dilayani. Sel semacam ini diberi biaya $M$ yang jauh lebih besar daripada biaya mana pun, sehingga metode akan menghindarinya selama masih ada alternatif (teknik Big-M). Di sini $M=10\times\max c_{ij}=3.512{,}1$ USD/kg. Bila di akhir sebuah solusi masih mengalokasikan barang pada sel ber-$M$, solusi itu tidak layak secara praktis dan biaya totalnya menjadi tidak bermakna. Selama iterasi MODI, sel ber-$M$ memiliki $d_{ij}$ yang sangat negatif sehingga cepat dikeluarkan dari basis.

\section{Data dan Metodologi}
\label{sec:metode}

\subsection{Sumber data}
Data berasal dari \textit{SCMS Delivery History Dataset} (versi 29 September 2015) yang dipublikasikan USAID \cite{usaid2015}. Berkas berisi 10.324 baris dan 33 kolom: identitas proyek, tanggal proses, nama vendor dan pabrik, deskripsi produk, harga, berat, dan biaya freight. Berkas tidak berenkode UTF-8 sehingga harus dibaca dengan enkode \texttt{latin-1}; bila tidak, karakter pada nama negara seperti \textit{C\^ote d'Ivoire} rusak. Data mencakup pengiriman ke 43 negara dari 88 nama lokasi pabrik berbeda.

\subsection{Alur penelitian}
Seluruh tahapan penelitian diringkas pada Gambar~\ref{fig:alur}. Data mentah dibersihkan terlebih dahulu (rincian pada Subbab~\ref{sec:cleaning}), kemudian diagregasi menjadi node sumber, node tujuan, pasokan, permintaan, dan matriks biaya. Matriks itu diselesaikan dengan tiga metode solusi awal yang berjalan sendiri-sendiri, dan tiap solusi awal diperbaiki dengan MODI. Hasil ketiganya lalu dibandingkan.

\begin{figure}[tbp]
\centering
\begin{tikzpicture}[node distance=0.42cm]
\node[ujung](s){Mulai};
\node[blok,below=of s](a){Mengambil dataset SCMS Delivery History\\(10.324 baris, 33 kolom)};
\node[blok,below=of a](b){Pembersihan data (Gambar~\ref{fig:clean})\\$\rightarrow$ 6.037 pengiriman, 4 kolom};
\node[blok,below=of b](c){Membentuk node sumber (22) dan tujuan (38);\\menghitung $a_i$, $b_j$, dan $c_{ij}$};
\node[blok,below=of c](d){Matriks biaya $22\times38$;\\sel kosong (578) diberi biaya Big-M};
\node[kecil,text width=2.4cm,below=0.6cm of d](lcm){LCM};
\node[kecil,text width=2.4cm,left=0.3cm of lcm](nwcr){NWCR};
\node[kecil,text width=2.4cm,right=0.3cm of lcm](vam){VAM};
\node[blok,below=0.6cm of lcm](m){Optimasi MODI untuk tiap solusi awal (Gambar~\ref{fig:modi})};
\node[blok,below=of m](h){Membandingkan biaya, jumlah iterasi, dan alokasi akhir};
\node[ujung,below=of h](e){Selesai};
\draw[panah](s)--(a); \draw[panah](a)--(b); \draw[panah](b)--(c); \draw[panah](c)--(d);
\draw[panah](d.south)--(lcm.north);
\draw[panah](d.south)--++(0,-0.3)-|(nwcr.north);
\draw[panah](d.south)--++(0,-0.3)-|(vam.north);
\draw[panah](nwcr.south)--(nwcr.south|-m.north);
\draw[panah](lcm.south)--(m.north);
\draw[panah](vam.south)--(vam.south|-m.north);
\draw[panah](m)--(h); \draw[panah](h)--(e);
\end{tikzpicture}
\caption{Alur penelitian dari pengambilan data sampai perbandingan hasil.}
\label{fig:alur}
\end{figure}

\subsection{Pembersihan data}
\label{sec:cleaning}
Kolom \texttt{Freight Cost (USD)} dan \texttt{Weight (Kilograms)} ternyata tidak seluruhnya berisi angka. Ada baris yang berisi teks seperti \textit{Freight Included in Commodity Cost}, \textit{Invoiced Separately}, \textit{Weight Captured Separately}, dan rujukan semacam \textit{See DN-304 (ID\#:10589)}. Rujukan itu menandakan bahwa biaya dan berat untuk satu pengiriman hanya ditulis sekali, di baris lain yang memiliki nomor ASN/DN sama. Bila angka-angka dijumlahkan per baris, biaya satu pengiriman terhitung berkali-kali. Pemeriksaan membenarkan hal ini: dari 7.030 nomor pengiriman unik, tidak ada yang memiliki lebih dari satu baris berbiaya numerik, dan tidak ada pengiriman yang menuju lebih dari satu negara. Oleh karena itu satuan analisis dipilih berupa pengiriman (ASN/DN), bukan baris item. Alur pembersihan ditunjukkan pada Gambar~\ref{fig:clean} dan rinciannya tercantum pada Tabel~\ref{tab:clean}.

\begin{figure}[tbp]
\centering
\begin{tikzpicture}[node distance=0.4cm]
\node[ujung](s){Dataset mentah: 10.324 baris, 33 kolom};
\node[blok,below=of s](a){Membaca CSV dengan enkode \texttt{latin-1}};
\node[blok,below=of a](b){Menghapus baris duplikat\\(semua kolom identik selain \texttt{ID})};
\node[samping,right=0.7cm of b](bb){4 baris dibuang\\sisa 10.320 baris};
\node[blok,below=of b](c){Mengubah freight dan berat menjadi angka (teks menjadi kosong); mengelompokkan baris per ASN/DN\\$\rightarrow$ 7.030 pengiriman};
\node[kep,below=of c](d){Ada baris dengan freight dan berat numerik $>0$?};
\node[samping,right=0.7cm of d](dd){856 pengiriman dibuang};
\node[blok,below=0.55cm of d](e){Mengambil biaya dan berat dari baris tersebut; menjumlahkan kuantitas dan nilai barang; sumber = pabrik bernilai terbesar\\$\rightarrow$ 6.174 pengiriman};
\node[blok,below=of e](f){Mengelompokkan 88 nama pabrik menjadi kelompok produsen; \textit{From RDC} menjadi node \texttt{SCMS RDC}};
\node[blok,below=of f](g){Menghitung biaya per kg = freight $\div$ berat};
\node[kep,below=of g](h){Biaya per kg di luar batas $3\times$IQR (skala log)?};
\node[samping,right=0.7cm of h](hh){19 pengiriman dibuang};
\node[kep,below=0.55cm of h](i){Sumber berkelompok ``Lainnya''?};
\node[samping,right=0.7cm of i](ii){118 pengiriman dibuang};
\node[blok,below=0.55cm of i](j){Menyimpan 4 kolom: \texttt{sumber}, \texttt{tujuan}, \texttt{berat\_kg}, \texttt{biaya\_freight\_usd}};
\node[ujung,below=of j](t){Data bersih: 6.037 pengiriman};
\draw[panah](s)--(a); \draw[panah](a)--(b); \draw[panah](b)--(c); \draw[panah](c)--(d);
\draw[panah](b)--(bb); \draw[panah](d)--(e) node[pos=0.4,right,font=\scriptsize]{ya};
\draw[panah](d)--(dd) node[pos=0.4,above,font=\scriptsize]{tidak};
\draw[panah](e)--(f); \draw[panah](f)--(g); \draw[panah](g)--(h);
\draw[panah](h)--(hh) node[pos=0.4,above,font=\scriptsize]{ya};
\draw[panah](h)--(i) node[pos=0.4,right,font=\scriptsize]{tidak};
\draw[panah](i)--(ii) node[pos=0.4,above,font=\scriptsize]{ya};
\draw[panah](i)--(j) node[pos=0.4,right,font=\scriptsize]{tidak};
\draw[panah](j)--(t);
\end{tikzpicture}
\caption{Alur pembersihan data. Kotak putus-putus menunjukkan data yang dibuang pada tiap tahap.}
\label{fig:clean}
\end{figure}

\begin{table}[tbp]
\centering\footnotesize
\caption{Tahapan pembersihan data.}\label{tab:clean}
\begin{tabular}{p{0.035\textwidth}p{0.6\textwidth}r}
\toprule
No & Perlakuan & Hasil \\ \midrule
1 & Membaca berkas CSV (enkode \texttt{latin-1}) & 10.324 baris, 33 kolom \\
2 & Menghapus duplikat (semua kolom identik selain \texttt{ID}); duplikat berdasarkan \texttt{ID} tidak ada & 4 baris dibuang, sisa 10.320 \\
3 & Memeriksa isi kolom biaya: 6.198 baris numerik; 1.442 \textit{included in commodity cost}, 239 \textit{invoiced separately}, 2.441 rujukan ke baris lain & 4.122 baris non-numerik \\
4 & Memeriksa isi kolom berat: 6.372 baris numerik; 1.507 \textit{captured separately}, 2.441 rujukan ke baris lain & 3.948 baris non-numerik \\
5 & Mengelompokkan baris menurut nomor ASN/DN & 7.030 pengiriman \\
6 & Mengambil pengiriman yang biaya dan beratnya numerik dan positif; kuantitas dan nilai barang dijumlahkan per pengiriman & 6.174 pengiriman \\
7 & Mengelompokkan 88 nama \texttt{Manufacturing Site} menurut perusahaan produsen; pengiriman \textit{From RDC} dijadikan satu node \texttt{SCMS RDC} & 21 produsen, 1 RDC, 1 ``Lainnya'' \\
8 & Menghitung biaya per kg dan membuang pencilan ekstrem pada skala log (di luar $Q_1-3\,\mathrm{IQR}$ sampai $Q_3+3\,\mathrm{IQR}$, atau 0,026--1.749 USD/kg) & 19 pengiriman dibuang \\
9 & Membuang pengiriman berkelompok ``Lainnya'' karena tidak mewakili satu lokasi produsen & 118 pengiriman dibuang \\
10 & Memilih kolom yang dibutuhkan: sumber, tujuan, berat, biaya freight & \textbf{6.037 pengiriman, 4 kolom} \\
\bottomrule
\end{tabular}
\end{table}

Beberapa catatan atas tahap-tahap tersebut. Pada tahap 7, nama seperti \textit{Cipla, Goa}, \textit{Cipla, Kurkumbh}, dan \textit{Cipla, Patalganga} dilebur menjadi satu node ``Cipla (India)'' karena berasal dari perusahaan yang sama; penggabungan serupa dilakukan untuk Aurobindo, Hetero, Mylan, dan produsen lain. Jika satu pengiriman berasal dari lebih dari satu pabrik (jarang terjadi), sumbernya ditetapkan sebagai pabrik dengan nilai barang terbesar dalam pengiriman itu. Pada tahap 8, contoh pencilan yang dibuang antara lain pengiriman seberat 1--2 kg dengan biaya freight ribuan dolar, dan pengiriman 857.354 kg dengan freight hanya 5.157 USD. Keduanya hampir pasti kesalahan pencatatan satuan atau angka, bukan tarif sungguhan. Dari 33 kolom awal, hanya empat yang dipertahankan sebagai input model (Tabel~\ref{tab:cols}); kolom lain seperti kode proyek, tanggal proses, nama vendor, deskripsi produk, harga satuan, dan asuransi tidak memengaruhi pasokan, permintaan, maupun biaya.

\begin{table}[tbp]
\centering\footnotesize
\caption{Kolom data yang dipakai pada model.}\label{tab:cols}
\begin{tabular}{lll}
\toprule
Kolom akhir & Kolom asal & Peran dalam model \\ \midrule
\texttt{sumber} & \texttt{Manufacturing Site}, \texttt{Fulfill Via} & node sumber $i$ \\
\texttt{tujuan} & \texttt{Country} & node tujuan $j$ \\
\texttt{berat\_kg} & \texttt{Weight (Kilograms)} & dasar $a_i$, $b_j$ \\
\texttt{biaya\_freight\_usd} & \texttt{Freight Cost (USD)} & dasar $c_{ij}$ \\
\bottomrule
\end{tabular}
\end{table}

\subsection{Pembentukan node, pasokan, permintaan, dan biaya}
Setiap node sumber adalah satu kelompok produsen atau gudang \texttt{SCMS RDC}, dan setiap node tujuan adalah satu negara. Notasi yang dipakai: $w_k$ untuk berat dan $f_k$ untuk biaya freight pengiriman ke-$k$, sedangkan $K_{ij}$ himpunan pengiriman dari $i$ ke $j$. Maka
\begin{equation}
a_i=\sum_{j}\sum_{k\in K_{ij}} w_k,\qquad
b_j=\sum_{i}\sum_{k\in K_{ij}} w_k,\qquad
c_{ij}=\frac{\sum_{k\in K_{ij}} f_k}{\sum_{k\in K_{ij}} w_k}.
\label{eq:param}
\end{equation}
Karena pasokan dan permintaan dihitung dari kumpulan pengiriman yang sama, jumlah keduanya otomatis identik, yaitu 19.550.180 kg, sehingga tidak diperlukan node fiktif (\textit{dummy}). Biaya rute dihitung sebagai rasio total freight terhadap total berat (rata-rata tertimbang) agar pengiriman besar berpengaruh lebih besar daripada pengiriman kecil. Hasilnya terbentuk $22+38=60$ node dan 258 rute teramati dari $836$ kemungkinan (30,9\%).

\subsection{Asumsi}
Pemodelan ini bertumpu pada asumsi-asumsi berikut.
\begin{enumerate}
\item Seluruh komoditas dianggap homogen dan diukur dalam kilogram, tanpa membedakan obat dan alat tes.
\item Ongkos angkut linear terhadap berat dan konstan pada setiap rute, yaitu sama dengan rata-rata tertimbang 2006--2015 tanpa koreksi inflasi.
\item Pasokan tiap sumber sama dengan volume yang tercatat pernah dikirim darinya, dan permintaan tiap negara sama dengan volume yang tercatat diterimanya.
\item Rute yang tidak pernah muncul di data dianggap tidak tersedia (biaya Big-M); tidak dilakukan estimasi biaya untuk rute tersebut.
\item Satu pengiriman dianggap berasal dari satu sumber utama, yaitu pabrik dengan nilai barang terbesar.
\item \texttt{SCMS RDC} diperlakukan sebagai sumber, bukan sebagai titik transit. Biaya yang dihitung hanya freight dari sumber ke tujuan, tanpa asuransi, bea, dan penanganan.
\item Moda angkut (udara, truk, laut) tidak dibedakan, dan tidak ada batas kapasitas pada rute.
\end{enumerate}

\subsection{Algoritma penyelesaian}
Alur ketiga metode solusi awal ditunjukkan pada Gambar~\ref{fig:awal}, dan alur iterasi MODI pada Gambar~\ref{fig:modi}. Tiga hal teknis perlu dicatat. Pertama, bila pasokan baris dan permintaan kolom habis bersamaan, hanya baris yang dicoret sehingga jumlah variabel dasar tetap $m+n-1$; kolom yang tertinggal menerima alokasi nol pada langkah berikutnya. Kedua, pada MODI nilai $u_i$ dan $v_j$ dihitung dengan menelusuri pohon sel basis mulai dari baris pertama, dan siklus dicari sebagai lintasan dari baris ke kolom sel masuk pada graf sel basis. Ketiga, iterasi berhenti ketika $d_{ij}\ge -10^{-9}$ untuk semua sel non-basis, dengan batas pengaman 5.000 iterasi yang tidak pernah tercapai.

\begin{figure}[tbp]
\centering
\begin{tikzpicture}[node distance=0.4cm]
\begin{scope}
\node[kecil](n1){\textbf{NWCR}\\Mulai dari sel $(1,1)$};
\node[kecil,below=of n1](n2){$x_{ij}=\min(a_i,b_j)$; kurangi $a_i$ dan $b_j$};
\node[kecil,below=of n2](n3){Jika $a_i=0$ pindah ke baris berikutnya; jika $b_j=0$ pindah ke kolom berikutnya};
\node[kep,text width=2.1cm,aspect=2,below=of n3](n4){Semua terpenuhi?};
\node[ujung,below=0.6cm of n4](n5){Solusi awal};
\draw[panah](n1)--(n2); \draw[panah](n2)--(n3); \draw[panah](n3)--(n4);
\draw[panah](n4)--(n5) node[pos=0.4,right,font=\scriptsize]{ya};
\draw[panah](n4.west)--([xshift=-0.3cm]n3.west|-n4.west)|-(n2.west) node[pos=0.2,above,font=\scriptsize]{tidak};
\end{scope}
\begin{scope}[xshift=5.05cm]
\node[kecil](l1){\textbf{LCM}\\Pilih sel aktif dengan $c_{ij}$ terkecil};
\node[kecil,below=of l1](l2){$x_{ij}=\min(a_i,b_j)$; kurangi $a_i$ dan $b_j$};
\node[kecil,below=of l2](l3){Coret baris atau kolom yang habis (bila keduanya habis, coret baris)};
\node[kep,text width=2.1cm,aspect=2,below=of l3](l4){Masih ada baris dan kolom aktif?};
\node[ujung,below=0.6cm of l4](l5){Solusi awal};
\draw[panah](l1)--(l2); \draw[panah](l2)--(l3); \draw[panah](l3)--(l4);
\draw[panah](l4)--(l5) node[pos=0.4,right,font=\scriptsize]{tidak};
\draw[panah](l4.west)--([xshift=-0.3cm]l3.west|-l4.west)|-(l1.west) node[pos=0.2,above,font=\scriptsize]{ya};
\end{scope}
\begin{scope}[xshift=10.1cm]
\node[kecil](v1){\textbf{VAM}\\Hitung penalti $p_i$ dan $q_j$ semua garis aktif};
\node[kecil,below=of v1](v2){Pilih garis dengan penalti terbesar};
\node[kecil,below=of v2](v3){Pilih sel termurah pada garis itu; $x_{ij}=\min(a_i,b_j)$};
\node[kecil,below=of v3](v4){Kurangi $a_i$ dan $b_j$; coret baris atau kolom yang habis};
\node[kep,text width=2.1cm,aspect=2,below=of v4](v5){Masih ada baris dan kolom aktif?};
\node[ujung,below=0.6cm of v5](v6){Solusi awal};
\draw[panah](v1)--(v2); \draw[panah](v2)--(v3); \draw[panah](v3)--(v4); \draw[panah](v4)--(v5);
\draw[panah](v5)--(v6) node[pos=0.4,right,font=\scriptsize]{tidak};
\draw[panah](v5.west)--([xshift=-0.3cm]v4.west|-v5.west)|-(v1.west) node[pos=0.2,above,font=\scriptsize]{ya};
\end{scope}
\end{tikzpicture}
\caption{Alur iterasi NWCR, LCM, dan VAM untuk memperoleh solusi dasar layak awal.}
\label{fig:awal}
\end{figure}

\begin{figure}[tbp]
\centering
\begin{tikzpicture}[node distance=0.4cm]
\node[ujung](s){Solusi dasar layak awal (NWCR, LCM, atau VAM)};
\node[blok,below=of s](a){Hitung $u_i$ dan $v_j$ dari sel basis: $u_i+v_j=c_{ij}$, dengan $u_1=0$};
\node[blok,below=of a](b){Hitung $d_{ij}=c_{ij}-u_i-v_j$ untuk semua sel non-basis};
\node[kep,below=of b](c){Semua $d_{ij}\ge 0$?};
\node[ujung,right=1.1cm of c](ok){Optimal: hitung $Z=\sum c_{ij}x_{ij}$};
\node[blok,below=0.55cm of c](d){Pilih sel masuk: $d_{ij}$ paling negatif};
\node[blok,below=of d](e){Cari siklus tertutup (sel masuk dan sel basis); beri tanda $+$ dan $-$ bergantian mulai dari sel masuk};
\node[blok,below=of e](f){$\theta=\min x_{ij}$ pada sel bertanda $-$};
\node[blok,below=of f](g){Tambah $\theta$ pada sel $+$, kurangi pada sel $-$; sel $-$ yang menjadi nol keluar dari basis};
\draw[panah](s)--(a); \draw[panah](a)--(b); \draw[panah](b)--(c);
\draw[panah](c)--(ok) node[pos=0.4,above,font=\scriptsize]{ya};
\draw[panah](c)--(d) node[pos=0.4,right,font=\scriptsize]{tidak};
\draw[panah](d)--(e); \draw[panah](e)--(f); \draw[panah](f)--(g);
\draw[panah](g.west)--++(-0.7,0)|-(a.west);
\end{tikzpicture}
\caption{Alur satu iterasi MODI sampai tercapai solusi optimal.}
\label{fig:modi}
\end{figure}

\subsection{Paket program}
Seluruh pengolahan dikerjakan dengan Python 3 pada \textit{Jupyter Notebook}, memakai \texttt{pandas} \cite{mckinney2010} untuk pembersihan dan agregasi, \texttt{NumPy} untuk komputasi matriks, serta \texttt{matplotlib} dan \texttt{seaborn} untuk grafik. NWCR, LCM, VAM, dan MODI ditulis sendiri sebagai fungsi terpisah dengan masukan matriks biaya, vektor pasokan, dan vektor permintaan. Tidak ada \textit{solver} program linear yang dipakai, sehingga seluruh iterasi dapat ditelusuri. Sebagai pemeriksaan, biaya akhir ketiga metode harus sama karena berasal dari masalah yang sama.

\section{Hasil dan Pembahasan}
\label{sec:hasil}

\subsection{Gambaran data}
Setelah pembersihan, tersisa 6.037 pengiriman dengan total berat 19.550.180 kg dan total freight 67.933.960 USD. Biaya rata-rata tertimbang adalah 3,475 USD/kg, sedangkan median biaya per pengiriman 7,12 USD/kg; selisih keduanya menunjukkan bahwa pengiriman besar cenderung lebih murah per kilogram. Berat pengiriman bersifat sangat miring ke kanan, dengan median 1.092 kg dan maksimum 112.027 kg. Moda udara mendominasi (3.965 pengiriman atau 65,7\%), disusul truk (19,2\%), \textit{air charter} (7,0\%), dan laut (4,7\%); sisanya 3,5\% tidak mencatat moda. Pengiriman via RDC mencakup 3.186 pengiriman (52,8\%) dan 70,2\% dari seluruh berat.

\begin{figure}[tbp]
\centering
\includegraphics[width=0.75\textwidth]{figs/f1_hist.pdf}
\caption{Sebaran $\log_{10}$ biaya per kg tiap pengiriman sebelum pencilan dibuang. Garis putus-putus menandai batas $Q_1-3\,\mathrm{IQR}$ dan $Q_3+3\,\mathrm{IQR}$.}
\label{fig:hist}
\end{figure}

Biaya per kg juga berubah menurut tahun (Gambar~\ref{fig:tahun}). Nilai tertimbang turun dari 7,27 USD/kg pada 2006 menjadi 2,44 USD/kg pada 2015, walaupun pada 2006 hanya ada 30 pengiriman sehingga angka tahun itu kurang kokoh. Penurunan ini sebagian mencerminkan perubahan komposisi muatan dan moda, bukan semata efisiensi. Karena model memakai seluruh tahun sekaligus, biaya rute yang dihasilkan merupakan nilai rata-rata sepanjang periode.

\begin{figure}[tbp]
\centering
\includegraphics[width=0.92\textwidth]{figs/f2_tahun.pdf}
\caption{Jumlah pengiriman dan biaya tertimbang per kg menurut tahun penerimaan.}
\label{fig:tahun}
\end{figure}

\subsection{Struktur node dan rute}
Pasokan sangat terpusat. \texttt{SCMS RDC} sendiri menyumbang 13.718.694 kg atau 70,17\% dari total, jauh di atas produsen terbesar berikutnya, Aurobindo, dengan 7,76\% (Tabel~\ref{tab:supply}, Gambar~\ref{fig:sd}). Permintaan lebih tersebar tetapi tetap timpang: sepuluh negara teratas (Tabel~\ref{tab:demand}) menyerap 90,7\% volume, dipimpin Nigeria dengan 3.866.619 kg, sementara enam negara lain menerima kurang dari seribu kilogram. Daftar permintaan lengkap ada pada Lampiran~\ref{app:demand}.

\begin{table}[tbp]
\centering\footnotesize
\caption{Node sumber dan pasokannya.}\label{tab:supply}
\begin{tabular}{rlrrr}\toprule No & Node sumber & Supply (kg) & Pangsa (\%) & Tujuan terhubung \\ \midrule
1 & SCMS RDC & 13.718.694 & 70,17 & 30 \\
2 & Aurobindo (India) & 1.517.691 & 7,76 & 21 \\
3 & Mylan/Matrix (India) & 763.471 & 3,91 & 17 \\
4 & Alere/Inverness (Japan) & 692.498 & 3,54 & 30 \\
5 & Hetero (India) & 654.646 & 3,35 & 16 \\
6 & Cipla (India) & 620.200 & 3,17 & 16 \\
7 & Strides (India) & 393.004 & 2,01 & 11 \\
8 & AbbVie/Abbott (Eropa) & 246.440 & 1,26 & 17 \\
9 & Trinity Biotech (Irlandia) & 210.178 & 1,08 & 20 \\
10 & KHB (China) & 209.605 & 1,07 & 6 \\
11 & Standard Diagnostics (Korea) & 134.591 & 0,69 & 12 \\
12 & GSK & 105.406 & 0,54 & 4 \\
13 & Chembio (USA) & 94.572 & 0,48 & 8 \\
14 & Micro Labs (India) & 58.706 & 0,30 & 9 \\
15 & Emcure (India) & 56.867 & 0,29 & 8 \\
16 & MSD & 23.718 & 0,12 & 5 \\
17 & BMS (Eropa) & 15.308 & 0,08 & 9 \\
18 & Ranbaxy (India) & 10.823 & 0,06 & 4 \\
19 & Novartis & 9.962 & 0,05 & 1 \\
20 & Orasure (USA) & 7.670 & 0,04 & 3 \\
21 & Roche & 5.257 & 0,03 & 8 \\
22 & Janssen & 873 & 0,00 & 3 \\
\midrule & \textbf{Total} & \textbf{19.550.180} & 100,00 & \\ \bottomrule\end{tabular}

\end{table}

\begin{figure}[tbp]
\centering
\includegraphics[width=0.95\textwidth]{figs/f3_sd.pdf}
\caption{Pasokan seluruh node sumber (skala log) dan permintaan 15 negara tujuan terbesar.}
\label{fig:sd}
\end{figure}

\begin{table}[tbp]
\centering\footnotesize
\caption{Sepuluh negara tujuan dengan permintaan terbesar.}\label{tab:demand}
\begin{tabular}{rlrrr}\toprule No & Negara tujuan & Demand (kg) & Pangsa (\%) & Sumber terhubung \\ \midrule
1 & Nigeria & 3.866.619 & 19,78 & 15 \\
2 & Zambia & 2.635.348 & 13,48 & 9 \\
3 & Mozambique & 2.347.865 & 12,01 & 9 \\
4 & South Africa & 1.697.753 & 8,68 & 6 \\
5 & Zimbabwe & 1.657.278 & 8,48 & 11 \\
6 & Côte d'Ivoire & 1.379.976 & 7,06 & 13 \\
7 & Tanzania & 1.368.242 & 7,00 & 15 \\
8 & Uganda & 1.257.919 & 6,43 & 17 \\
9 & Rwanda & 882.445 & 4,51 & 14 \\
10 & Vietnam & 635.320 & 3,25 & 12 \\
\bottomrule\end{tabular}

\end{table}

Matriks biaya pada Gambar~\ref{fig:heat} memperlihatkan bahwa jaringan ini jarang. Dari 836 sel, hanya 258 yang terisi dan sisanya, 578 sel, tidak pernah dipakai. Dari 258 rute itu, 44 hanya didukung satu pengiriman sehingga biayanya rentan terhadap kasus tunggal. Biaya per kg rute berkisar 0,51--351,21 USD/kg dengan median 8,67 dan rata-rata (tak tertimbang) 15,37. Hanya \texttt{SCMS RDC} dan Alere/Inverness yang terhubung ke 30 negara; produsen kecil seperti Novartis hanya melayani satu negara. Kondisi inilah yang membuat sel kosong harus diperlakukan hati-hati.

\begin{figure}[tbp]
\centering
\includegraphics[width=\textwidth]{figs/f4_heat.pdf}
\caption{Matriks biaya $c_{ij}$ (USD/kg, skala log). Sel putih adalah rute yang tidak pernah muncul di data.}
\label{fig:heat}
\end{figure}

\subsection{Solusi awal NWCR, LCM, dan VAM}
Ketiga metode dijalankan pada matriks biaya dengan Big-M, dan semuanya menghasilkan 59 sel dasar sesuai teori. Perbedaannya terletak pada kelayakan dan biaya (Tabel~\ref{tab:awal}).

\begin{table}[tbp]
\centering\footnotesize\setlength{\tabcolsep}{3pt}
\caption{Solusi awal ketiga metode. Biaya rute valid hanya menghitung alokasi pada rute yang pernah teramati.}\label{tab:awal}
\begin{tabular}{lrrrrl}
\toprule
Metode & Sel pada & Berat pada & Biaya rute valid & Biaya dengan & Status \\
 & rute terlarang & rute terlarang & (USD) & Big-M (USD) & \\ \midrule
NWCR & 37 & 2.413.372 kg (12,3\%) & 72.073.871 & 8.548.077.672 & tidak layak \\
LCM & 8 & 275.944 kg (1,4\%) & 62.022.909 & 1.031.165.832 & tidak layak \\
VAM & 0 & 0 & 56.691.684 & 56.691.684 & layak \\
\bottomrule
\end{tabular}
\end{table}

\textbf{NWCR.} Metode ini membagi barang menurut posisi sel, sehingga sebagian besar pembagian jatuh pada pasangan yang tidak pernah ada. Dari 59 sel dasar, 37 berada pada rute terlarang dengan total 2.413.372 kg (12,3\% dari seluruh volume). Dengan memasukkan Big-M, biayanya mencapai sekitar 8,5 miliar USD, angka yang hanya menegaskan bahwa solusinya tidak dapat dipakai apa adanya.

\textbf{LCM.} Metode ini jauh lebih baik karena sel murah diambil lebih dulu, tetapi tidak sempurna. Masih ada 8 sel dasar pada rute terlarang yang memuat 275.944 kg (1,4\%). Hal ini terjadi karena LCM menghabiskan sel-sel murah lebih dulu, sehingga tersisa pasangan sumber--tujuan yang hanya bisa dipertemukan lewat rute yang tidak ada. Biaya total dengan Big-M mendekati 1,03 miliar USD.

\textbf{VAM.} Seluruh 59 sel dasar VAM jatuh pada rute yang valid, sehingga solusinya layak dengan biaya 56.691.684 USD, atau 16,5\% di bawah biaya aktual (67.933.960 USD).

\subsection{Hasil optimasi MODI}
Setiap solusi awal kemudian dimasukkan ke MODI (Tabel~\ref{tab:modi}). Ketiganya berhenti pada biaya yang sama, yaitu 54.461.845 USD, tanpa satu pun alokasi pada rute terlarang. Jumlah iterasi yang dibutuhkan berbeda jauh: 97 untuk NWCR, 42 untuk LCM, dan 16 untuk VAM. Gambar~\ref{fig:konv} memperlihatkan penurunan biaya pada tiap iterasi; garis NWCR dan LCM mulai dari biaya Big-M yang sangat tinggi lalu jatuh ke wilayah biaya aktual dalam beberapa iterasi pertama, sedangkan VAM sudah berada di dekat optimum sejak awal.

\begin{table}[tbp]
\centering\footnotesize\setlength{\tabcolsep}{3pt}
\caption{Hasil MODI untuk tiap solusi awal.}\label{tab:modi}
\begin{tabular}{lrrrr}
\toprule
Metode & Biaya awal (USD) & Iterasi MODI & Biaya akhir (USD) & Sel terlarang akhir \\ \midrule
NWCR & 8.548.077.672 & 97 & 54.461.845 & 0 \\
LCM & 1.031.165.832 & 42 & 54.461.845 & 0 \\
VAM & 56.691.684 & 16 & 54.461.845 & 0 \\
\bottomrule
\end{tabular}
\end{table}

\begin{figure}[tbp]
\centering
\includegraphics[width=0.85\textwidth]{figs/f5_konv.pdf}
\caption{Biaya total (skala log, termasuk Big-M) pada tiap iterasi MODI. Garis titik-titik menandai biaya aktual.}
\label{fig:konv}
\end{figure}

Alokasi akhir ketiga metode identik sel demi sel (Lampiran~\ref{app:final}). Biaya reduksi terkecil pada sel non-basis bernilai 0,14 USD/kg, positif, sehingga solusi optimal ini tunggal. Seluruh 59 sel basis berisi alokasi positif, dan 15 sel di antaranya menggantikan sel pada solusi awal VAM. Contohnya, kiriman \texttt{SCMS RDC} ke Mozambique turun dari 2.347.865 kg menjadi 1.946.462 kg, sedangkan Alere/Inverness yang semula tidak melayani Mozambique kini mengirim 401.403 kg. Di Afrika Selatan, bagian Aurobindo naik dari 526.438 kg menjadi 1.419.467 kg, sementara bagian Cipla turun dari 620.200 kg menjadi 278.286 kg. Sebanyak 15 tujuan dilayani langsung oleh \texttt{SCMS RDC} pada solusi akhir, dengan total 13.718.694 kg, sama dengan pasokannya.

\subsection{Pembahasan}
Perbedaan jumlah iterasi sejalan dengan kualitas solusi awal. NWCR tidak melihat biaya sama sekali, sedangkan data ini memuat 69\% sel terlarang sehingga alokasi acak hampir pasti menabraknya; MODI harus mengeluarkan 37 sel ber-$M$ satu per satu dan baru kemudian menyempurnakan biaya. LCM melihat biaya tetapi hanya secara lokal, sehingga sisa 8 sel terlarang dan 42 iterasi. VAM mempertimbangkan seberapa mahal kerugian bila sebuah sel murah tidak dipakai. Pada matriks jarang, baris atau kolom yang hanya punya satu atau dua rute valid mendapat penalti sangat besar (karena sel lainnya bernilai $M$) sehingga langsung diprioritaskan. Itulah sebabnya VAM praktis terhindar dari rute terlarang dan hanya memerlukan 16 iterasi.

Dari sisi biaya, solusi optimal lebih murah 13.472.115 USD (19,83\%) daripada biaya aktual. Solusi awal VAM sudah 16,5\% di bawah biaya aktual dan hanya 4,09\% di atas optimum; MODI memangkas sisa selisih 2.229.839 USD. Karena ketiga metode berakhir pada solusi yang sama dan optimum bersifat tunggal, pilihan metode awal tidak memengaruhi hasil akhir. Pilihan itu hanya menentukan beban hitung, dan di sini VAM jelas paling hemat.

Hasil ini perlu dibaca hati-hati. Pertama, penghematan dihitung dari biaya rata-rata per rute, padahal tarif aktual dipengaruhi moda, ukuran muatan, dan tahun. Kedua, sumber di dunia nyata tidak dapat saling menggantikan: obat tertentu hanya dibuat oleh produsen tertentu, sehingga pasokan satu pabrik tidak bisa dialihkan ke negara mana pun. Ketiga, sebagian rute hanya didukung satu pengiriman. Jadi angka penghematan lebih tepat dipandang sebagai batas atas potensi, bukan prediksi yang pasti tercapai.

\section{Kesimpulan dan Saran}
\label{sec:kesimpulan}

\subsection{Kesimpulan}
\begin{enumerate}
\item Data SCMS dapat dijadikan masalah transportasi setelah dibersihkan dengan tepat. Kunci utamanya adalah memakai pengiriman (ASN/DN) sebagai satuan analisis agar biaya tidak terhitung ganda. Dari 10.324 baris mentah diperoleh 6.037 pengiriman dengan empat kolom, yang membentuk 22 node sumber, 38 node tujuan, dan 258 rute teramati.
\item Dengan pasokan dan permintaan dari total berat pengiriman, masalahnya berimbang pada 19.550.180 kg. Sebanyak 578 dari 836 sel tidak tersedia sehingga perlu dilarang dengan Big-M.
\item NWCR dan LCM menghasilkan solusi awal yang tidak layak pada data ini (masing-masing 12,3\% dan 1,4\% berat berada pada rute terlarang), sedangkan VAM menghasilkan solusi awal layak dengan biaya 56.691.684 USD.
\item Setelah dioptimalkan dengan MODI, ketiga metode mencapai biaya optimal yang sama, 54.461.845 USD, atau 19,83\% di bawah biaya aktual. Jumlah iterasi MODI sebesar 97 (NWCR), 42 (LCM), dan 16 (VAM) menunjukkan bahwa VAM adalah titik awal yang paling efisien untuk data ini.
\end{enumerate}

\subsection{Saran}
Biaya untuk rute kosong dapat diperkirakan, misalnya lewat regresi terhadap jarak dan moda, supaya tidak seluruhnya dilarang. Peran RDC lebih tepat dimodelkan sebagai titik transit (\textit{transshipment}), dan analisis akan lebih tajam bila dipisah per kelompok produk serta dibatasi pada tahun-tahun terbaru, misalnya 2012--2015, agar tarif lebih seragam. Kapasitas rute dan kendala pasokan per produk perlu dimasukkan agar hasilnya lebih dekat dengan kondisi operasional. Bagi pengelola data, pencatatan freight yang masih berupa teks (\textit{included in commodity cost}, \textit{invoiced separately}) layak dibenahi karena membuat sekitar sepertiga pengiriman tidak bisa dianalisis.

\appendix
\section{Solusi optimal hasil MODI}
\label{app:final}
Tabel~\ref{tab:final} memuat 59 sel basis solusi optimal. Hasil yang sama diperoleh dari solusi awal NWCR, LCM, maupun VAM.
\begingroup\footnotesize\setlength{\tabcolsep}{4pt}
\begin{longtable}{rllrrr}\caption{Solusi optimal hasil MODI (59 sel basis, identik untuk solusi awal NWCR, LCM, dan VAM)}\label{tab:final}\\\toprule No & Sumber & Tujuan & Alokasi (kg) & $c_{ij}$ (USD/kg) & Biaya (USD) \\ \midrule\endfirsthead
\toprule No & Sumber & Tujuan & Alokasi (kg) & $c_{ij}$ (USD/kg) & Biaya (USD) \\ \midrule\endhead
1 & SCMS RDC & Nigeria & 3.856.657 & 2,68 & 10.326.970 \\
2 & SCMS RDC & Zambia & 2.635.348 & 1,65 & 4.347.534 \\
3 & SCMS RDC & Mozambique & 1.946.462 & 0,89 & 1.741.305 \\
4 & SCMS RDC & Zimbabwe & 1.526.030 & 1,67 & 2.542.671 \\
5 & SCMS RDC & Tanzania & 1.362.985 & 2,07 & 2.819.607 \\
6 & SCMS RDC & Côte d'Ivoire & 1.255.014 & 3,36 & 4.213.584 \\
7 & SCMS RDC & Rwanda & 823.760 & 4,29 & 3.535.331 \\
8 & SCMS RDC & Ghana & 119.313 & 2,44 & 290.563 \\
9 & SCMS RDC & Namibia & 72.936 & 4,30 & 313.618 \\
10 & SCMS RDC & Congo, DRC & 59.902 & 8,96 & 536.920 \\
11 & SCMS RDC & Burundi & 28.714 & 5,85 & 167.899 \\
12 & SCMS RDC & South Sudan & 17.754 & 10,80 & 191.823 \\
13 & SCMS RDC & Benin & 12.015 & 4,57 & 54.942 \\
14 & SCMS RDC & Libya & 1.783 & 4,69 & 8.361 \\
15 & SCMS RDC & Senegal & 21 & 30,13 & 633 \\
16 & Aurobindo (India) & South Africa & 1.419.467 & 1,47 & 2.082.074 \\
17 & Aurobindo (India) & Swaziland & 63.752 & 2,79 & 177.849 \\
18 & Aurobindo (India) & Ethiopia & 34.472 & 3,73 & 128.560 \\
19 & Mylan/Matrix (India) & Uganda & 706.291 & 3,29 & 2.321.720 \\
20 & Mylan/Matrix (India) & Malawi & 57.180 & 1,42 & 81.116 \\
21 & Alere/Inverness (Japan) & Mozambique & 401.403 & 8,62 & 3.461.900 \\
22 & Alere/Inverness (Japan) & Haiti & 254.718 & 10,33 & 2.632.358 \\
23 & Alere/Inverness (Japan) & Botswana & 25.047 & 2,06 & 51.607 \\
24 & Alere/Inverness (Japan) & Pakistan & 7.998 & 12,99 & 103.864 \\
25 & Alere/Inverness (Japan) & Sierra Leone & 1.719 & 18,45 & 31.710 \\
26 & Alere/Inverness (Japan) & Lesotho & 856 & 12,93 & 11.064 \\
27 & Alere/Inverness (Japan) & Mali & 725 & 34,07 & 24.704 \\
28 & Alere/Inverness (Japan) & Guinea & 32 & 56,38 & 1.804 \\
29 & Hetero (India) & Vietnam & 578.453 & 2,67 & 1.546.147 \\
30 & Hetero (India) & Haiti & 73.868 & 2,31 & 170.310 \\
31 & Hetero (India) & Togo & 2.325 & 5,06 & 11.754 \\
32 & Cipla (India) & South Africa & 278.286 & 1,38 & 382.699 \\
33 & Cipla (India) & Haiti & 230.484 & 3,09 & 712.311 \\
34 & Cipla (India) & Uganda & 111.430 & 3,88 & 431.914 \\
35 & Strides (India) & Uganda & 393.004 & 3,56 & 1.400.784 \\
36 & AbbVie/Abbott (Eropa) & Cameroon & 196.671 & 6,96 & 1.368.653 \\
37 & AbbVie/Abbott (Eropa) & Guyana & 24.758 & 7,71 & 190.924 \\
38 & AbbVie/Abbott (Eropa) & Dominican Republic & 23.097 & 2,49 & 57.454 \\
39 & AbbVie/Abbott (Eropa) & Angola & 1.914 & 8,70 & 16.656 \\
40 & Trinity Biotech (Irlandia) & Kenya & 168.433 & 5,87 & 988.382 \\
41 & Trinity Biotech (Irlandia) & Uganda & 31.013 & 8,64 & 268.033 \\
42 & Trinity Biotech (Irlandia) & Guyana & 10.732 & 10,84 & 116.352 \\
43 & KHB (China) & Ethiopia & 209.605 & 9,77 & 2.047.799 \\
44 & Standard Diagnostics (Korea) & Zimbabwe & 131.248 & 7,33 & 962.652 \\
45 & Standard Diagnostics (Korea) & Sudan & 2.676 & 12,50 & 33.440 \\
46 & Standard Diagnostics (Korea) & Liberia & 667 & 25,77 & 17.187 \\
47 & GSK & Côte d'Ivoire & 101.244 & 2,64 & 267.132 \\
48 & GSK & Guatemala & 4.162 & 7,22 & 30.048 \\
49 & Chembio (USA) & Ethiopia & 94.572 & 5,30 & 501.052 \\
50 & Micro Labs (India) & Rwanda & 58.685 & 2,39 & 140.240 \\
51 & Micro Labs (India) & Afghanistan & 21 & 25,03 & 526 \\
52 & Emcure (India) & Vietnam & 56.867 & 3,07 & 174.604 \\
53 & MSD & Côte d'Ivoire & 23.718 & 2,31 & 54.841 \\
54 & BMS (Eropa) & Uganda & 15.308 & 0,51 & 7.775 \\
55 & Ranbaxy (India) & Cameroon & 10.823 & 6,98 & 75.546 \\
56 & Novartis & Nigeria & 9.962 & 9,29 & 92.508 \\
57 & Orasure (USA) & Cameroon & 7.670 & 19,70 & 151.130 \\
58 & Roche & Tanzania & 5.257 & 6,49 & 34.132 \\
59 & Janssen & Uganda & 873 & 7,76 & 6.773 \\
\midrule & \multicolumn{2}{l}{\textbf{Total}} & \textbf{19.550.180} & & \textbf{54.461.845} \\ \bottomrule\end{longtable}

\endgroup

\section{Permintaan seluruh negara tujuan}
\label{app:demand}
\begingroup\footnotesize\setlength{\tabcolsep}{4pt}
\begin{longtable}{rlrrr}\caption{Permintaan seluruh 38 negara tujuan}\label{tab:demandfull}\\\toprule No & Negara tujuan & Demand (kg) & Pangsa (\%) & Sumber terhubung \\ \midrule\endfirsthead
\toprule No & Negara tujuan & Demand (kg) & Pangsa (\%) & Sumber terhubung \\ \midrule\endhead
1 & Nigeria & 3.866.619 & 19,78 & 15 \\
2 & Zambia & 2.635.348 & 13,48 & 9 \\
3 & Mozambique & 2.347.865 & 12,01 & 9 \\
4 & South Africa & 1.697.753 & 8,68 & 6 \\
5 & Zimbabwe & 1.657.278 & 8,48 & 11 \\
6 & Côte d'Ivoire & 1.379.976 & 7,06 & 13 \\
7 & Tanzania & 1.368.242 & 7,00 & 15 \\
8 & Uganda & 1.257.919 & 6,43 & 17 \\
9 & Rwanda & 882.445 & 4,51 & 14 \\
10 & Vietnam & 635.320 & 3,25 & 12 \\
11 & Haiti & 559.070 & 2,86 & 17 \\
12 & Ethiopia & 338.649 & 1,73 & 9 \\
13 & Cameroon & 215.164 & 1,10 & 10 \\
14 & Kenya & 168.433 & 0,86 & 5 \\
15 & Ghana & 119.313 & 0,61 & 7 \\
16 & Namibia & 72.936 & 0,37 & 8 \\
17 & Swaziland & 63.752 & 0,33 & 5 \\
18 & Congo, DRC & 59.902 & 0,31 & 8 \\
19 & Malawi & 57.180 & 0,29 & 3 \\
20 & Guyana & 35.490 & 0,18 & 10 \\
21 & Burundi & 28.714 & 0,15 & 4 \\
22 & Botswana & 25.047 & 0,13 & 7 \\
23 & Dominican Republic & 23.097 & 0,12 & 8 \\
24 & South Sudan & 17.754 & 0,09 & 5 \\
25 & Benin & 12.015 & 0,06 & 4 \\
26 & Pakistan & 7.998 & 0,04 & 2 \\
27 & Guatemala & 4.162 & 0,02 & 3 \\
28 & Sudan & 2.676 & 0,01 & 3 \\
29 & Togo & 2.325 & 0,01 & 2 \\
30 & Angola & 1.914 & 0,01 & 5 \\
31 & Libya & 1.783 & 0,01 & 1 \\
32 & Sierra Leone & 1.719 & 0,01 & 2 \\
33 & Lesotho & 856 & 0,00 & 1 \\
34 & Mali & 725 & 0,00 & 2 \\
35 & Liberia & 667 & 0,00 & 3 \\
36 & Guinea & 32 & 0,00 & 1 \\
37 & Afghanistan & 21 & 0,00 & 1 \\
38 & Senegal & 21 & 0,00 & 1 \\
\midrule & \textbf{Total} & \textbf{19.550.180} & 100,00 & \\ \bottomrule\end{longtable}

\endgroup

\bibliographystyle{unsrtnat}
\bibliography{jhep-indo}

\end{document}
